% Options for packages loaded elsewhere
\PassOptionsToPackage{unicode}{hyperref}
\PassOptionsToPackage{hyphens}{url}
\PassOptionsToPackage{dvipsnames,svgnames,x11names}{xcolor}
%
\documentclass[
  11pt,
  letterpaper,
]{article}
\usepackage{amsmath,amssymb}
\usepackage{iftex}
\ifPDFTeX
  \usepackage[T1]{fontenc}
  \usepackage[utf8]{inputenc}
  \usepackage{textcomp} % provide euro and other symbols
\else % if luatex or xetex
  \usepackage{unicode-math} % this also loads fontspec
  \defaultfontfeatures{Scale=MatchLowercase}
  \defaultfontfeatures[\rmfamily]{Ligatures=TeX,Scale=1}
\fi
\usepackage{lmodern}
\ifPDFTeX\else
  % xetex/luatex font selection
  \setmainfont[]{DejaVu Serif}
  \setsansfont[]{DejaVu Sans}
  \setmonofont[]{DejaVu Sans Mono}
\fi
% Use upquote if available, for straight quotes in verbatim environments
\IfFileExists{upquote.sty}{\usepackage{upquote}}{}
\IfFileExists{microtype.sty}{% use microtype if available
  \usepackage[]{microtype}
  \UseMicrotypeSet[protrusion]{basicmath} % disable protrusion for tt fonts
}{}
\makeatletter
\@ifundefined{KOMAClassName}{% if non-KOMA class
  \IfFileExists{parskip.sty}{%
    \usepackage{parskip}
  }{% else
    \setlength{\parindent}{0pt}
    \setlength{\parskip}{6pt plus 2pt minus 1pt}}
}{% if KOMA class
  \KOMAoptions{parskip=half}}
\makeatother
\usepackage{xcolor}
\usepackage[margin=0.8in]{geometry}
\usepackage{longtable,booktabs,array}
\usepackage{calc} % for calculating minipage widths
% Correct order of tables after \paragraph or \subparagraph
\usepackage{etoolbox}
\makeatletter
\patchcmd\longtable{\par}{\if@noskipsec\mbox{}\fi\par}{}{}
\makeatother
% Allow footnotes in longtable head/foot
\IfFileExists{footnotehyper.sty}{\usepackage{footnotehyper}}{\usepackage{footnote}}
\makesavenoteenv{longtable}
\usepackage{graphicx}
\makeatletter
\def\maxwidth{\ifdim\Gin@nat@width>\linewidth\linewidth\else\Gin@nat@width\fi}
\def\maxheight{\ifdim\Gin@nat@height>\textheight\textheight\else\Gin@nat@height\fi}
\makeatother
% Scale images if necessary, so that they will not overflow the page
% margins by default, and it is still possible to overwrite the defaults
% using explicit options in \includegraphics[width, height, ...]{}
\setkeys{Gin}{width=\maxwidth,height=\maxheight,keepaspectratio}
% Set default figure placement to htbp
\makeatletter
\def\fps@figure{htbp}
\makeatother
\setlength{\emergencystretch}{3em} % prevent overfull lines
\providecommand{\tightlist}{%
  \setlength{\itemsep}{0pt}\setlength{\parskip}{0pt}}
\setcounter{secnumdepth}{-\maxdimen} % remove section numbering

\usepackage{microtype}
\usepackage{needspace}
\usepackage{seqsplit}
\usepackage{caption}
\usepackage{float}
\makeatletter\def\fps@figure{H}\makeatother
\renewcommand{\figurename}{Figura}
\renewcommand{\contentsname}{Conteúdo}
\captionsetup{font=small,labelfont=bf}
\setlength{\emergencystretch}{3em}
\setlength{\parskip}{0.45em}
\setlength{\parindent}{0pt}
\setcounter{tocdepth}{1}
\widowpenalty=10000
\clubpenalty=10000
\ifLuaTeX
  \usepackage{selnolig}  % disable illegal ligatures
\fi
\IfFileExists{bookmark.sty}{\usepackage{bookmark}}{\usepackage{hyperref}}
\IfFileExists{xurl.sty}{\usepackage{xurl}}{} % add URL line breaks if available
\urlstyle{same}
\hypersetup{
  pdftitle={Escolha do regime operacional e limites de informação no controle inferencial de reatores químicos},
  colorlinks=true,
  linkcolor={black},
  filecolor={Maroon},
  citecolor={Blue},
  urlcolor={black},
  pdfcreator={LaTeX via pandoc}}

\title{Escolha do regime operacional e limites de informação no controle
inferencial de reatores químicos}
\author{}
\date{10 de outubro de 2026}

\begin{document}
\maketitle

Estudo computacional exploratório de propilenoglicol e etilbenzeno

10 de outubro de 2026

\hypertarget{resumo}{%
\section{Resumo}\label{resumo}}

Investiga-se se a operação pode facilitar o controle de composição sem
medição química contínua, preservando dois modelos históricos de
reatores. A campanha combina auditoria física, excitação térmica, ciclos
com recursos comparáveis, mapas de sensibilidade e políticas congeladas
sob incerteza. No propilenoglicol, excitação senoidal reduziu a
incerteza de uma qualidade futura, sem remover dois modos químicos
invisíveis. A seleção do regime reduziu o erro médio do mesmo PI térmico
em 41,8\% com qualidade nominal igual, mas diminuiu a produção em 2,24\%
e apresentou violações térmicas. No etilbenzeno, a melhora de tracking
desapareceu quando a qualidade foi igualada. Nenhum dos ciclos
admissíveis superou um estacionário encontrado com recursos iguais ou
menores. Uma ablação termodinâmica mostrou que um pequeno defeito de
Hess pode aumentar artificialmente o posto de observabilidade local. Os
resultados sustentam uma metodologia de falsificação de ganhos
inferenciais, sem estabelecer dominância operacional, segurança
industrial ou um princípio universal.

Palavras chave: controle inferencial; operação de reatores;
observabilidade; invariantes; projeto de experimentos.

\hypertarget{introduuxe7uxe3o-e-revisuxe3o-cruxedtica}{%
\section{Introdução e revisão
crítica}\label{introduuxe7uxe3o-e-revisuxe3o-cruxedtica}}

O desempenho de um sistema de controle depende tanto da política quanto
do regime físico em que ela atua. Quando a composição não é medida
continuamente, mudar a alimentação ou a temperatura pode alterar
simultaneamente a produtividade, a informação térmica e a sensibilidade
da qualidade. A pergunta relevante é se essa intervenção melhora
decisões sob a mesma especificação e restrições.

Operação periódica não é uma ideia nova, inclusive para propilenoglicol
{[}R01{]}. Resposta não linear em frequência e manipulação com defasagem
já foram usadas para avaliar sistemas reacionais {[}R03, R04{]}.
Resultados recentes em metanol mostram atividade atual nessa linha
{[}R06-R08{]}, mas não autorizam transferir benefícios para outras
cinéticas. Controle auto-otimizante relaciona a escolha de variáveis à
perda econômica {[}R13{]}. Projeto de experimento voltado à decisão
explicita que precisão paramétrica e utilidade operacional são objetivos
diferentes {[}R20, R21{]}.

Há antecedentes igualmente importantes para os resultados negativos.
Invariantes impõem restrições de observabilidade em CSTRs {[}R09,
R10{]}; avanços de reação podem simplificar modelos e controladores
{[}R11, R22{]}. A modelagem por entalpia interfere no projeto de
observadores {[}R12{]}, e propriedades termodinâmicas podem exigir
reconciliação {[}R23{]}. A contribuição pretendida aqui é testar essas
relações nos benchmarks recuperados, expondo comparações confundidas e
limites de interpretação. Não se afirma novidade internacional de cada
componente.

\hypertarget{modelos-fuxedsicos-e-proceduxeancia}{%
\section{Modelos físicos e
procedência}\label{modelos-fuxedsicos-e-proceduxeancia}}

O PG parte da dissertação de Brandão de 2019 {[}R25{]} e dos códigos
históricos e reconstruídos associados. Os seis estados são as
concentrações A, B, produto C e inerte M, além das temperaturas do
reator e do trocador. A cinética histórica é de primeira ordem em A. Com
volume fixo, taxa r e diluição D:

\[
\dot c=D(c_f-c)+\nu r,\qquad
\nu=(-1,-1,1,0)^T,\qquad r=k(T)C_A.
\]

O balanço térmico usa capacidade proporcional à soma das capacidades
molares ponderadas pelas concentrações. O calor do trocador depende
diretamente da vazão refrigerante e da temperatura do reator; a
temperatura dinâmica do trocador não realimenta o balanço do reator.
Essa assimetria foi preservada, em vez de substituída por uma camisa
convencional. O limite principal é 160 °F, equivalente a 344,261 K; o
limite posterior de 180 °F não foi usado para aceitar trajetórias.

O EB é a reconstrução do primeiro reator do TFC de Brandão {[}R26{]},
associado ao processo discutido por Luyben {[}R14{]}. Seus estados são
E, B, EB, DEB e temperatura. As reações são E+B para EB, E+EB para DEB e
B+DEB para dois EB. As três taxas foram preservadas, assim como o
reciclo de entrada nominal. O volume é 200 m³. O produto controlado é a
fração molar de EB; a produção líquida desconta o EB que entra no
reciclo. O fluxograma com separação e reciclos dinâmicos não foi
executado.

A reconstrução EB não é uma planta termodinamicamente reconciliada. Os
coeficientes históricos tratados como capacidades caloríficas coincidem
aproximadamente com entalpia dividida por temperatura em °C. O dever
nominal da reconstrução é 25,370 MW, enquanto o backup Aspen registra
9,088 MW em uma condição próxima, mas distinta. O envelope de 140,85 a
165 °C é exploratório, não industrial certificado. Os novos ensaios
assumem acesso a um dever térmico idealizado; essa autoridade adicional
é declarada.

\hypertarget{formulauxe7uxe3o-cientuxedfica}{%
\section{Formulação científica}\label{formulauxe7uxe3o-cientuxedfica}}

Considere a incerteza em inventários e parâmetros, representada por uma
distribuição b compatível com o histórico de entradas e temperaturas. A
escolha de regime z e política causal π deve minimizar custo
operacional, incluindo transições, qualidade, produção, recursos e
risco. Uma medida relevante de falta de informação é o arrependimento:

\[
\mathcal R(b,z)=\min_\pi {\mathbb E}_b[J(\xi,z,\pi)]
-{\mathbb E}_b[\min_\pi J(\xi,z,\pi)].
\]

A matriz de Fisher fornece um instrumento local, mas não define esse
valor. Menor variância de qualidade pode resultar de menor incerteza,
menor sensibilidade ou menor influência do inventário inicial. Para uma
classe de condições indistinguíveis, uma garantia comum exige interseção
não vazia entre suas políticas admissíveis que atendem à especificação.
Mudar o regime pode ampliar essa interseção ou eliminá-la por restrições
térmicas. A campanha aproxima esse problema por mapas, políticas fixas e
decisões finitas; não resolve a otimização global sobre todas as
políticas.

No PG, duas perturbações químicas com A e temperaturas iguais e
capacidade térmica total igual permanecem termicamente indistinguíveis.
Sua diferença segue

\[
\delta c(t)=\delta c(0)\exp\left[-\int_0^t D(s)\,ds\right].
\]

Excitação comum não identifica esses modos. Uma separação inicial de
produto de 0,5 kmol/m³ impede garantia simultânea de uma banda de ±2\%
antes de pelo menos 0,240 h, mesmo usando o maior D permitido. Essa é
uma condição necessária; limites de velocidade podem torná-la mais
restritiva. O antecedente são os invariantes, e não uma propriedade
criada pelo estimador.

\hypertarget{muxe9todos-experimentais}{%
\section{Métodos experimentais}\label{muxe9todos-experimentais}}

A triagem avaliou 44 projetos de excitação, 144 ciclos e 376 linhas de
paisagem estacionária. Os sinais incluem degrau, pulso realizável, PRBS,
chirp, senoides e multissenoides. As derivadas incluem inventário
inicial e parâmetros térmicos desconhecidos, com escalas físicas
mantidas entre regiões. Os ciclos foram avaliados após fechamento
periódico, com verificação da partida e médias reais dos atuadores.

O teste não linear de informação compreendeu 192 ajustes MAP, quatro
sinais e 24 sementes por reator, das quais oito continham bias fora da
família ajustada. A comparação decisória principal usa as 16 sementes da
família correta e 17 ações estacionárias futuras comuns, com oráculo de
parâmetros claramente identificado. O custo inclui qualidade e
penalização térmica; não representa lucro nem controle ótimo dinâmico.

Na malha fechada, os ganhos de um PI térmico com filtro e integral foram
selecionados no nominal e congelados. A política não recebe composição
verdadeira nem usa a composição estimada em paralelo. A campanha inicial
de regiões contém 2304 execuções rotuladas e famílias de mismatch
cinético, térmico, combinado, instrumental e fora do domínio. Regiões
coincidentes não são réplicas independentes.

Após detectar que a comparação inicial mudava a referência química,
executou-se uma rodada exploratória adicional de 216 trajetórias, com
três famílias e 12 sementes inéditas. Ajustou-se a alimentação para
manter exatamente a mesma qualidade nominal nos pontos comparados. Os
recursos resultantes foram reportados, sem presumir igualdade.
Intervalos de 95\% usam bootstrap pareado de 4000 reamostragens por
semente, sem tratar as famílias da mesma semente como independentes. Não
houve correção de multiplicidade.

Os baselines adaptados incluem C0, C2 e C3 no PG e reconstrução FKE+PI,
cascata e NMPC com envelope no EB, totalizando 180 trajetórias. Ensaios
adicionais cobrem esquecimento, amostra atrasada, supervisão, gêmeos e
coordenadas reduzidas. Quinze testes científicos e o fluxo E2E foram
executados. Dez trajetórias foram reintegradas por DOP853, com diferença
térmica máxima de 0,000214 K em relação à integração principal.

\hypertarget{resultados}{%
\section{Resultados}\label{resultados}}

\hypertarget{informauxe7uxe3o-nuxe3o-equivale-a-decisuxe3o}{%
\subsection{Informação não equivale a
decisão}\label{informauxe7uxe3o-nuxe3o-equivale-a-decisuxe3o}}

No PG, a melhor senoide reduziu o desvio posterior local da qualidade
futura de 0,053965 para 0,008416 kmol/m³. Na validação não linear, o
RMSE caiu de 0,016685 para 0,006019. O arrependimento médio passou de
0,105895 a 0,004239; a diferença teve IC95\% de -0,183124 a -0,034132.
Entretanto, o degrau teve arrependimento médio 0,002931, menor que o da
senoide. Logo, a ordenação por precisão preditiva não coincidiu com a
ordenação decisória.

No EB, a redução local de desvio posterior foi 7,1\%. O arrependimento
médio da senoide caiu de 1,496401 para 0,811008, mas o IC da diferença,
de -1,695583 a 0,190828, incluiu zero. Não se demonstrou separação
confiável das três contribuições cinéticas nem um projeto universalmente
superior.

\hypertarget{periodicidade-e-seleuxe7uxe3o-de-regime}{%
\subsection{Periodicidade e seleção de
regime}\label{periodicidade-e-seleuxe7uxe3o-de-regime}}

Foram admissíveis 68 ciclos PG e 45 EB. Os maiores ganhos de produção
sobre o nominal foram apenas 0,164\% e 0,031\%. Frente a estacionários
encontrados com recursos iguais ou menores, os melhores ciclos ficaram
12,622\% e 0,898\% abaixo, respectivamente. Esses resultados refutam a
dominância da família examinada, não toda forma de operação periódica.

Os mapas identificaram regiões produtivas menos sensíveis em ambos os
modelos. No PG, o ponto econômico substituto e o de menor sensibilidade
diferiram. No EB, quatro critérios selecionaram o mesmo ponto para os
pesos iniciais, mas a seleção mudou sob outros pesos. Menor
sensibilidade não foi confundida com maior posto de observabilidade.

Na rodada com qualidade comum, os resultados do PI congelado foram:

\begin{longtable}[]{@{}lrrr@{}}
\toprule\noalign{}
Comparação com nominal & PG a 337 K & PG a 340 K & EB a 162 °C \\
\midrule\noalign{}
\endhead
\bottomrule\noalign{}
\endlastfoot
Variação do IAE normalizado & -30,85\% & -41,82\% & -0,26\% \\
Variação da produção & -1,59\% & -2,24\% & -0,12\% \\
Variação da alimentação & -8,33\% & -11,98\% & -0,72\% \\
Trajetórias com violação & 4/36 & 3/36 & 0/36 \\
\end{longtable}

No PG a 340 K, a diferença de IAE foi -0,014114, IC95\% {[}-0,022466;
-0,007626{]}. A produção caiu 0,638944 kmol/h, IC95\% {[}-0,886668;
-0,348401{]}. O nominal também violou limites em 3/36 trajetórias; a
duração média das violações aumentou na alternativa. A condição não é
aceita como segura. No EB, o IC da diferença de IAE incluiu zero e o
ganho praticamente desapareceu.

\begin{figure}
\centering
\includegraphics[width=0.95\textwidth,height=\textheight]{figures/same_quality.png}
\caption{Efeitos com a mesma qualidade e política congelada. Intervalos
exploratórios pareados por semente.}
\end{figure}

\hypertarget{informauxe7uxe3o-aparente-por-inconsistuxeancia-fuxedsica}{%
\subsection{Informação aparente por inconsistência
física}\label{informauxe7uxe3o-aparente-por-inconsistuxeancia-fuxedsica}}

O fechamento isolado do ciclo de Hess trouxe uma falsificação adicional.
Seja N a matriz estequiométrica, R a derivada das taxas e D a diluição.
O bloco químico linearizado é \(A_c=-DI+NR\). Se \(\Delta H^T=h^TN\),
então todo \(v\in\ker(NR)\) satisfaz \(A_cv=-Dv\) e \(\Delta H^TRv=0\).
Como os demais termos térmicos da reconstrução não dependem da
concentração, esses modos ficam invisíveis à temperatura na
linearização.

Nos doze equilíbrios examinados, o núcleo tem dimensão dois. O residual
histórico de Hess, 5,5 kcal/kmol, acrescentava uma direção térmica
espúria, elevando o posto local de três para quatro. No nominal, o
quarto valor singular do registro normalizado passou de 0,001183 para
aproximadamente 9e-14 após o fechamento. Seu autovalor de Fisher
histórico era apenas 1,4e-6: posto maior não significava informação
prática apreciável.

\begin{figure}
\centering
\includegraphics[width=0.9\textwidth,height=\textheight]{figures/hess_observability.png}
\caption{A direção adicional desaparece ao fechar Hess. A ablação não
reconcilia a planta EB.}
\end{figure}

Outro contrafactual trocou incorretamente a referência de temperatura e
reajustou Q no nominal. O equilíbrio pôde ser preservado, mas o ganho
térmico de alimentação mudou de 1,515 para 4,138 °C por (m³/h). Uma
calibração escalar estacionária, portanto, pode esconder uma derivada
física incorreta.

\hypertarget{discussuxe3o-e-limitauxe7uxf5es}{%
\section{Discussão e limitações}\label{discussuxe3o-e-limitauxe7uxf5es}}

Os resultados favorecem a seleção de regime como variável de pesquisa,
mas refutam uma interpretação universalmente favorável. O PG permite
reduzir erro sem trocar controlador; ainda assim, os cenários térmicos
invalidam a promoção da alternativa. O EB evidencia o risco de atribuir
melhoria à informação quando parte do efeito vem de mudança de
referência. Transferiu-se uma filosofia experimental, não parâmetros ou
uma política ótima.

Três falsificações mudaram a interpretação. Um piloto periódico parecia
ganhar aproximadamente 5,4\% de produção PG porque a limitação de
velocidade deslocava a média da alimentação; a correção eliminou o
ganho. A renovação acumulada igual produziu a mesma ambiguidade terminal
dos gêmeos PG, apesar de diferenças transitórias. A ablação de Hess
mostrou que até o posto formal pode ser favorecido por uma
inconsistência.

O atlas adversarial é finito e não certifica um núcleo de viabilidade.
Os oráculos de estado e parâmetros são restritos às políticas avaliadas,
não ótimos globais. A amostra química é um protocolo separado, com
atraso causal; sua seleção por variância não dominou horários fixos em
controle. O supervisor e as combinações não demonstraram ganho adicional
seguro.

Toda a evidência nova é computacional. MATLAB, Aspen e planta física não
foram executados. A termodinâmica EB, a idealização de remoção de calor,
o reciclo fixo, os priors sintéticos, o pequeno número de sementes
independentes e a busca finita limitam inferências industriais. A
revisão é dirigida e encerrada em 10/10/2026. O efeito local de Hess tem
antecedentes em invariantes e observadores por entalpia e requer revisão
adicional antes de uma reivindicação de originalidade.

\hypertarget{conclusuxf5es}{%
\section{Conclusões}\label{conclusuxf5es}}

Modificar a operação pode melhorar inferência ou reduzir a informação
necessária para algumas decisões, mas os mecanismos e os custos precisam
ser separados. Nos modelos estudados, não foi demonstrada uma região
segura que melhorasse simultaneamente qualidade, produção e informação
nos dois processos. A contribuição mais madura é um protocolo
reprodutível que confronta ganhos aparentes com recursos iguais,
qualidade comum, limites informacionais e consistência física. Os
resultados negativos e a observabilidade espúria merecem investigação
adicional, sem serem apresentados como uma nova tecnologia industrial.

\hypertarget{proceduxeancia-e-disponibilidade}{%
\section{Procedência e
disponibilidade}\label{proceduxeancia-e-disponibilidade}}

O código, dados brutos, sementes, modelos, testes e figuras acompanham o
pacote da campanha. Os modelos históricos são de Walter Yanko de Aragão
Brandão {[}R25, R26{]}. A participação documentada de Arioston Araújo de
Morais Júnior é de orientação no TFC e banca/agradecimentos na
dissertação; a orientação do mestrado é de Heleno Bispo da Silva Júnior.
O manuscrito não atribui nova coautoria, contribuição individual ou
afiliação. Autoria e destino de submissão exigem revisão dos
pesquisadores.

\hypertarget{referuxeancias}{%
\section{Referências}\label{referuxeancias}}

R01. Pérez-Polo, M. F.; Pérez-Molina, M. Increasing the reactant
conversion through induced oscillations in a continuous stirred tank
reactor by using PI control. \emph{Journal of Process Control}, 23,
778-792, 2013. DOI
\href{https://doi.org/10.1016/j.jprocont.2013.03.007}{10.1016/j.jprocont.2013.03.007}.
Verificação: resumo do editor.

R03. Nikolić, D.; Seidel-Morgenstern, A.; Petkovska, M. Periodic
Operation with Modulation of Inlet Concentration and Flow Rate. Part I:
Nonisothermal CSTR. \emph{Chemical Engineering \& Technology}, 39,
2020-2028, 2016. DOI
\href{https://doi.org/10.1002/ceat.201600185}{10.1002/ceat.201600185}.
Verificação: editor.

R04. Živković, L. A.; Vidaković-Koch, T.; Petkovska, M. Computer-Aided
Nonlinear Frequency Response Method for Investigating the Dynamics of
Chemical Engineering Systems. \emph{Processes}, 8, 1354, 2020. DOI
\href{https://doi.org/10.3390/pr8111354}{10.3390/pr8111354}. Texto
integral consultado pelo repositório PSE/LAPSE.

R06. Leipold, J.; Nikolić, D.; Seidel-Morgenstern, A.; Kienle, A.
Optimization of methanol synthesis under forced periodic operation in a
non-isothermal fixed-bed reactor. \emph{Computers \& Chemical
Engineering}, 2025, 109040. DOI
\href{https://doi.org/10.1016/j.compchemeng.2025.109040}{10.1016/j.compchemeng.2025.109040}.
Verificação: editor e Max Planck.

R07. Kaps, L. et al.~Forced Periodic Reactor Operation Applied to
Methanol Synthesis. \emph{ChemCatChem}, 2026, e202501403. DOI
\href{https://doi.org/10.1002/cctc.202501403}{10.1002/cctc.202501403}.
Primeiro acesso online publicado em 14/12/2025; verificação: editor.

R08. Leipold, J.; Kienle, A. Forced periodic operation of chemical
reactors: Origins of performance improvement demonstrated for methanol
synthesis. \emph{Chemical Engineering Science}, publicação online em
2026, 124639. DOI
\href{https://doi.org/10.1016/j.ces.2026.124639}{10.1016/j.ces.2026.124639}.
Verificação: resumo do editor e instituição; não confundir ano online
com ano do fascículo.

R09. Fjeld, M.; Asbjørnsen, O. A.; Åström, K. J. Reaction invariants and
their importance in the analysis of eigenvectors, state observability
and controllability of the continuous stirred tank reactor.
\emph{Chemical Engineering Science}, 29, 1917-1926, 1974. DOI
\href{https://doi.org/10.1016/0009-2509(74)85009-8}{10.1016/0009-2509(74)85009-8}.
Verificação: editor e Universidade de Lund.

R10. Srinivasan, B.; Amrhein, M.; Bonvin, D. Reaction and flow
variants/invariants in chemical reaction systems with inlet and outlet
streams. \emph{AIChE Journal}, 44, 1858-1867, 1998. DOI
\href{https://doi.org/10.1002/aic.690440815}{10.1002/aic.690440815}.
Verificação: editor.

R11. Amrhein, M.; Bhatt, N.; Srinivasan, B.; Bonvin, D. Extents of
reaction and flow for homogeneous reaction systems with inlet and outlet
streams. \emph{AIChE Journal}, 2010. DOI
\href{https://doi.org/10.1002/aic.12125}{10.1002/aic.12125}.
Verificação: editor.

R12. Dochain, D.; Couenne, F.; Jallut, C. Enthalpy based modelling and
design of asymptotic observers for chemical reactors.
\emph{International Journal of Control}, 82, 1389-1403, 2009. DOI
\href{https://doi.org/10.1080/00207170802236085}{10.1080/00207170802236085}.
Verificação: resumo do editor; acesso integral não obtido.

R13. Skogestad, S. Plantwide control: the search for the self-optimizing
control structure. \emph{Journal of Process Control}, 10, 487-507, 2000.
DOI
\href{https://doi.org/10.1016/S0959-1524(00)00023-8}{10.1016/S0959-1524(00)00023-8}.
Texto integral do manuscrito do autor consultado na NTNU.

R14. Luyben, W. L. Design and control of the ethylbenzene process.
\emph{AIChE Journal}, 2011. DOI
\href{https://doi.org/10.1002/aic.12289}{10.1002/aic.12289}.
Verificação: editor e PDF no acervo TFC. O fluxograma completo não é
intercambiável com o R1 isolado.

R20. Wan, B.; Liu, M.; Grigas, P.; Shen, Z.-J. M. Decision-Focused
Sequential Experimental Design: A Directional Uncertainty-Guided
Approach. Preprint de 05/02/2026.
\href{https://arxiv.org/abs/2602.05340}{arXiv:2602.05340}. Não tratado
como artigo revisado por pares nem como validação química.

R21. Go, J.; Qian, X.; Yoon, B.-J. Goal-driven Bayesian Optimal
Experimental Design for Robust Decision-Making Under Model Uncertainty.
Preprint de 25/05/2026.
\href{https://arxiv.org/abs/2605.26093}{arXiv:2605.26093}. Verificação:
arXiv; aplicação não constitui demonstração nos reatores PG/EB.

R22. Marquez Ruiz, A.; Mendez Blanco, C.; Özkan, L. Constrained Control
and Estimation of Homogeneous Reaction Systems Using Extent-Based Linear
Parameter Varying Models. \emph{Industrial \& Engineering Chemistry
Research}, 59, 2242-2251, 2020. DOI
\href{https://doi.org/10.1021/acs.iecr.9b04412}{10.1021/acs.iecr.9b04412}.
Verificação: editor e registro dos autores na TU Eindhoven.

R23. Salike, S.; Bhatt, N. Thermodynamically consistent estimation of
Gibbs free energy from data: data reconciliation approach.
\emph{Bioinformatics}, 36, 1219-1225, 2020; online em 04/10/2019. DOI
\href{https://doi.org/10.1093/bioinformatics/btz741}{10.1093/bioinformatics/btz741}.
Texto integral consultado. A reconciliação de Gibbs é antecedente
metodológico; não fornece os Cp faltantes do EB.

R25. Brandão, W. Y. de A. \emph{Uso de observadores de estado aplicados
no controle inferencial preditivo de processos não lineares}.
Dissertação, UFCG, 2019. Fonte primária recuperada no acervo do usuário;
orientação de Heleno Bispo da Silva Júnior, participação de Arioston
Araújo de Morais Júnior na banca/agradecimentos.

R26. Brandão, W. Y. de A. \emph{Uso de sensores virtuais como
observadores de estado aplicados na dinâmica e controle de processos
químicos lineares e não lineares}. Trabalho final de curso, UFPB, 2016.
Fonte primária recuperada; orientação de Arioston Araújo de Morais
Júnior.

\end{document}
